The article \cite{SalaMartinDoppelhoferMiller_BACE} suggests a
technique for averaging the parameter estimates -- Bayesian Averaging
of Classical Estimates, which enables to measure the importance of
particular potential regressors. This approach is alternative to the
familiar and earlier-applied Bayesian Model Averaging (BMA) technique,
from which it differs mainly because of its less restrictive \textit{a
  priori} assumptions. The authors showed that a weighting method may
be understood as a limiting case of Bayesian analysis when the prior
information is dominated by the data. Averaging of the parameter
estimates is performed across all possible combinations of models
obtained by means of Classical Ordinary Least Squares (OLS)
estimation.

The BACE approach has several interesting features in relation to
the techniques such as BMA or robustness-checking methods.
First of all, it is not necessary to make assumptions concerning
the \textit{a priori} distribution for parameters, and only
one simple assumption, on the expected model size.
% The weights for particular models
% are built based on the logarithm of the likelihood function
% corrected for degrees of freedom. 
% The estimation of parameters
% is realised by means of Classical OLS. Moreover, there are
% considered all possible models and there are not taken any
% assumptions concerning the individual regressors (no variable is of
% fixed type).

% Let's start with Bayes' rule:
% 
% \begin{equation}\label{eq:Bayes rule}
%   g(\beta | y)= \frac{f(y|\beta )g(\beta )}{f(y)}.
% \end{equation}
% 
% In equation (\ref{eq:Bayes rule}) $g(\beta)$ is the prior density
% of vector $\beta$ and $f(y)$ is density of observed data. Given
% two subsets of variables (two models) – $M_0$ and $M_1$ – 
% under Bayes' rule we get:
% 
% \begin{equation}\label{eq:refText6}
%     g(\beta | y) = P(M_0)\frac{f(y|\beta )g(\beta | M_0)}{f(y)}
%     + P(M_1)\frac{f(y|\beta)g(\beta|M_1)}{f(y)}.
% \end{equation}
% Now, we can rewrite (\ref{eq:refText6}) in terms of posterior probabilities of two models:
% 
% \begin{equation}
%     g(\beta |y) = P(M_0|y)\frac{f(y|\beta )g(\beta |M_0)}{f(y|M_0)}
%     + P(M_1|y)\frac{f(y|\beta )g(\beta |M_1)}{f(y|M_1)},
% \end{equation}
% where $P(M_i|y)$ is the posterior probability of the i'th model.

By applying diffuse priors we obtain the posterior odds ratio for two
regression models:
\begin{equation}
\frac{P(M_0|y)}{P(M_1|y)} \approx
\frac{P(M_0)}{P(M_1)}\left(\frac{|A_0|/|A_0 + X_0^{'}X_0|}{|A_1|/|A_1 + X_1^{'}X_1|}\right)^{-1/2}
\left(\frac{SSE_0+Q_0}{SSE_1+Q_1}\right)^{-T/2},
\end{equation}
where
\begin{itemize}
\item $P(M_i)$ and $P(M_i|y)$ are the prior and posterior
  probabilities, respectively, for model $M_i$;
\item $X_0$ and $X_1$ are matrices with the observations for the
  independent variables for $M_0$ and $M_1$;
  $SSE_i$ is the OLS sum of squared errors for model $i$;
\item $T$ is the sample size;
\item $A_0^{-1}$ and $A_1^{-1}$ are covariance matrices;
\item  $Q_i$ is a quadratic form of estimated parameters.
\end{itemize}

If we use Zellner's \textit{g}-prior \cite{Zellner1986} for
$A_0=\mathit{gX_0}^{'}X_0$ and $A_1=\mathit{gX_1}^{'}X_1$
and assume that $g\rightarrow 0$ and $X^{T}X\rightarrow \infty$
we will finally obtain the following posterior odds ratio:
\begin{equation}\label{eq:refText9}
\frac{P(M_0|y)}{P(M_1|y)}=\frac{P(M_0)}{P(M_1)}T^{(k_1-k_0)/2}
\left(\frac{SSE_0}{SSE_1}\right)^{-T/2},
\end{equation}
where $k_i$ is the number of parameters in model $M_i$.  We can
calculate posterior probability of $P(M_j|y)$ via the following
formula:

\begin{equation}
P(M_j|y)=\frac{P(M_j)T^{-k_j/2}SSE_j^{-T/2}}
    {\sum_{i=1}^{2^K}P(M_i)T^{-k_i/2}SSE_i^{-T/2}}.
\end{equation}
where $2^K$ denotes the total number of potential combinations of the
\textit{K} independent variables.

Using BACE we can also easily evaluate the mean and variance of
the posterior distribution of slope parameters (see \cite{Leamer1978}):
\begin{equation}\label{eq:refText11}
E(\beta |y)=\sum_{i=1}^{2^K}P(M_j|y)E(\beta |y,M_j)
\end{equation}

Equation (\ref{eq:refText11}) can be rewritten as:
\begin{equation}
E(\beta |y)=\sum _{i=1}^{2^K}P(M_j|y)\hat{\beta}
\end{equation}
where $\hat{\beta}=E(\beta |y,M_j)$ is the OLS estimate of $\beta$
from model $M_j$. The posterior variance of $\beta$ is given by:
\begin{equation}
Var(\beta |y) = 
\sum_{i=1}^{2^K}P(M_j|y)Var(\beta |y,M_j) +
\sum_{i=1}^{2^K}P(M_j|y)\left(\hat{\beta}_j -
\sum_{i=1}^{2^K}P(M_j|y)\hat{\beta}_j\right)^2.
\end{equation}

An example of the application of the BACE technique for forecasting
macroeconomic indicators can be found, among others, in the articles
\cite{vanDijk2004Forecasting,BialowolskiKuszewskiWitkowski2015Empirica,
  AlbisMapa2014MPRA55902}.

%%% Local Variables:
%%% mode: latex
%%% TeX-master: "BACE_package_for_gretl"
%%% End:
